\section{Existence and topological properties of Gibbs measures}

Throughout, $S$ is countable, $(E,\mathcal E)$ is a measurable space, $\Cfg = E^S$, and
$\gamma$ is a specification. Topological statements refer to the topology $\mathcal L$ of
local convergence on $\Prob(\Cfg,\mathcal F)$ unless the weak topology is named.

\subsection*{Finite-volume Gibbs distributions and the fixed-point equation}

\begin{definition}[Finite-volume Gibbs distribution $\nu\gamma_\Lambda$]
    \label{def:exg-finite-volume-distribution}
    \uses{def:spec, def:kernel}
    \lean{Specification.bindPM, MeasureTheory.GibbsMeasure.finiteVolumeDistributions}
    \leanok

    For $\nu\in\Prob(\Cfg,\mathcal F)$ and $\Lambda\in\Fin$,
    \[\nu\gamma_\Lambda(A) \;=\; \int \gamma_\Lambda(A\mid\eta)\,\nu(d\eta)\]
    is a probability measure. For $\eta\in\Cfg$, the net
    $\Lambda\mapsto\gamma_\Lambda(\cdot\mid\eta) = \delta_\eta\gamma_\Lambda$ indexed by
    $\Fin$ is the net of \textbf{finite-volume Gibbs distributions with boundary condition
    $\eta$}.
\end{definition}

\begin{lemma}[Gibbs measures as fixed points]
    \label{lem:exg-mem-gp-iff}
    \uses{def:gibbs-meas, def:exg-finite-volume-distribution}
    \lean{MeasureTheory.GibbsMeasure.mem_GP_iff_forall_bindPM_eq}
    \leanok

    For $\mu\in\Prob(\Cfg,\mathcal F)$, $\mu\in\Gibbs{\gamma}$ if and only if
    $\mu\gamma_\Lambda = \mu$ for all $\Lambda\in\Fin$.
\end{lemma}
\begin{proof}
    \leanok

    Since $\gamma_\Lambda$ is proper and $\mathcal F_{S\setminus\Lambda}$-measurable,
    $\mu\gamma_\Lambda = \mu$ is the DLR equation
    $\mu(A\mid\mathcal F_{S\setminus\Lambda}) = \gamma_\Lambda(A\mid\cdot)$ $\mu$-a.s.,
    $A\in\mathcal F$.
\end{proof}

\begin{definition}[Local thermodynamic limit, Comment (4.18)]
    \label{def:exg-local-thermo-limit}
    \uses{def:exg-finite-volume-distribution}
    \lean{MeasureTheory.GibbsMeasure.IsLocalThermodynamicLimit}
    \leanok

    $\mu\in\Prob(\Cfg,\mathcal F)$ is a \textbf{local thermodynamic limit} for the boundary
    condition $\eta$ if $\mu$ is an $\mathcal L$-cluster point of the net
    $(\gamma_\Lambda(\cdot\mid\eta))_{\Lambda\in\Fin}$.
\end{definition}

\subsection*{Georgii (4.17), (4.18), (4.22): existence}

\begin{theorem}[Georgii (4.17)]
    \label{thm:exg-4-17}
    \uses{def:gibbs-meas, def:exg-finite-volume-distribution, lem:exg-mem-gp-iff}
    \lean{MeasureTheory.GibbsMeasure.mem_GP_of_tendsto_withLocalConvergence,
      MeasureTheory.GibbsMeasure.mem_GP_of_mapClusterPt}
    \leanok

    Let $\gamma$ be quasilocal, $(\gamma^i)_{i}$ a net of specifications with
    $\gamma^i\to\gamma$ uniformly on local observables, $(\Lambda_i)_i$ a net in $\Fin$ with
    $\Lambda_i\to S$, and $(\nu_i)_i$ a net in $\Prob(\Cfg,\mathcal F)$. If
    $\nu_i\gamma^i_{\Lambda_i}\to\mu$ in $\mathcal L$, then $\mu\in\Gibbs{\gamma}$. The same
    holds for every cluster point $\mu$ of $(\nu_i\gamma^i_{\Lambda_i})_i$.
\end{theorem}
\begin{proof}
    \uses{lem:exg-mem-gp-iff}
    \leanok

    Fix $\Lambda\in\Fin$ and a local event $A$. The observable $g=\gamma_\Lambda \mathbf 1_A$
    is quasilocal, hence $\mathcal L$-continuous, so $\int g\,d(\nu_i\gamma^i_{\Lambda_i})
    \to\int g\,d\mu$. Uniform convergence replaces $g$ by $\gamma^i_\Lambda\mathbf 1_A$ with
    error tending to $0$, and for $\Lambda\subseteq\Lambda_i$ consistency gives
    $\int\gamma^i_\Lambda\mathbf 1_A\,d(\nu_i\gamma^i_{\Lambda_i})
    = (\nu_i\gamma^i_{\Lambda_i})(A)\to\mu(A)$. Hence $\mu\gamma_\Lambda=\mu$ on local
    events, which separate probability measures. For cluster points pass to an ultrafilter
    refining the index filter.
\end{proof}

\begin{corollary}[Georgii, Comment (4.18)]
    \label{cor:exg-4-18}
    \uses{thm:exg-4-17, def:exg-local-thermo-limit}
    \lean{MeasureTheory.GibbsMeasure.IsLocalThermodynamicLimit.mem_GP}
    \leanok

    If $\gamma$ is quasilocal, every local thermodynamic limit is a Gibbs measure for
    $\gamma$.
\end{corollary}
\begin{proof}
    \uses{thm:exg-4-17}
    \leanok

    Theorem \ref{thm:exg-4-17} with $\gamma^i=\gamma$, $\nu_i=\delta_\eta$, $\Lambda_i=\Lambda$,
    and $\delta_\eta\gamma_\Lambda = \gamma_\Lambda(\cdot\mid\eta)$.
\end{proof}

\begin{theorem}[Georgii (4.22)]
    \label{thm:exg-4-22}
    \uses{thm:exg-4-17, def:exg-local-thermo-limit}
    \lean{MeasureTheory.GibbsMeasure.exists_mem_GP_mapClusterPt,
      MeasureTheory.GibbsMeasure.exists_isLocalThermodynamicLimit_mem_GP}
    \leanok

    Let $(E,\mathcal E)$ be standard Borel and $\gamma$ quasilocal. If
    $\mu_i=\nu_i\gamma^i_{\Lambda_i}$ is a locally equicontinuous net with
    $\gamma^i\to\gamma$ uniformly on local observables and $\Lambda_i\to S$, then
    $\Gibbs{\gamma}$ contains a cluster point of $(\mu_i)_i$; in particular
    $\Gibbs{\gamma}\neq\emptyset$. For $\gamma^i\equiv\gamma$ and $\nu_i=\delta_\eta$ the Gibbs
    measure obtained is a local thermodynamic limit for $\eta$.
\end{theorem}
\begin{proof}
    \uses{thm:exg-4-17}
    \leanok

    By Georgii (4.9) a locally equicontinuous net over a standard Borel state space converges
    along any ultrafilter refining the index filter to some $\mu\in\Prob(\Cfg,\mathcal F)$;
    Theorem \ref{thm:exg-4-17} along that ultrafilter gives $\mu\in\Gibbs{\gamma}$.
\end{proof}

\subsection*{Topological properties of $\Gibbs{\gamma}$}

\begin{lemma}[$\Gibbs{\gamma}$ is $\mathcal L$-closed]
    \label{lem:exg-gp-closed}
    \uses{thm:exg-4-17, def:gibbs-meas}
    \lean{MeasureTheory.GibbsMeasure.isClosed_setOf_mem_GP}
    \leanok

    If $\gamma$ is quasilocal, $\Gibbs{\gamma}$ is $\mathcal L$-closed.
\end{lemma}
\begin{proof}
    \uses{thm:exg-4-17, lem:exg-mem-gp-iff}
    \leanok

    Let $\mu$ be a cluster point of $\Gibbs{\gamma}$. Index by
    $\bar D = (\mathcal N(\mu)\cap\Gibbs{\gamma})\times\Fin$ with
    $\gamma^{(\nu,\Lambda)}=\gamma$, $\nu_{(\nu,\Lambda)}=\nu$,
    $\Lambda_{(\nu,\Lambda)}=\Lambda$. Then $\nu\gamma_\Lambda=\nu\to\mu$, and
    Theorem \ref{thm:exg-4-17} gives $\mu\in\Gibbs{\gamma}$.
\end{proof}

\begin{lemma}[Uniform kernel bounds pass to Gibbs measures]
    \label{lem:exg-gp-dominated}
    \uses{lem:exg-mem-gp-iff}
    \lean{MeasureTheory.GibbsMeasure.apply_le_of_mem_GP}
    \leanok

    If $\mu\in\Gibbs{\gamma}$, $A\in\mathcal F$, $\Lambda\in\Fin$ and
    $\gamma_\Lambda(A\mid\omega)\le c$ for all $\omega$, then $\mu(A)\le c$.
\end{lemma}
\begin{proof}
    \uses{lem:exg-mem-gp-iff}
    \leanok

    $\mu(A)=(\mu\gamma_\Lambda)(A)=\int\gamma_\Lambda(A\mid\omega)\,\mu(d\omega)\le c$.
\end{proof}

\begin{example}[Georgii, Example (4.11)(1)]
    \label{ex:exg-4-11-1}
    \uses{def:isssd-spec, def:modif, lem:exg-gp-dominated, lem:exg-gp-closed}
    \lean{MeasureTheory.GibbsMeasure.isCompact_closure_setOf_mem_GP_modification}
    \leanok

    Let $(E,\mathcal E)$ be standard Borel, $\nu$ a single-spin probability measure and
    $\rho$ a modifier of $\ISSSD(\nu)$ with $\rho_\Lambda\le C_\Lambda<\infty$ for every
    $\Lambda\in\Fin$. Then $\Gibbs{\rho\,\ISSSD(\nu)}$ is relatively $\mathcal L$-compact.
\end{example}
\begin{proof}
    \uses{lem:exg-gp-dominated}
    \leanok

    On $\mathcal F_\Lambda$-events the kernel is dominated by $C_\Lambda\,\nu^S$, so by
    Lemma \ref{lem:exg-gp-dominated} every Gibbs measure lies in the set of random fields
    dominated by the finite measures $C_\Lambda\,\nu^S$, which is $\mathcal L$-compact by
    Georgii (4.9)--(4.11).
\end{proof}

\subsection*{Georgii (2.11): the space $\Bspace$ of absolutely summable potentials}

\begin{definition}[Georgii (2.11)--(2.12): $\Bspace$ and the seminorms $\|\cdot\|_i$]
    \label{def:exg-b-space}
    \lean{Potential.absolutelySummable, Potential.seminormAt, Potential.seminormFamily,
      Potential.withSeminorms_seminormFamily}
    \leanok

    $\Bspace$ is the $\R$-submodule of potentials $\Phi$ with $\Phi_\emptyset=0$, $\Phi_A$
    $\mathcal F_A$-measurable, and
    \[\|\Phi\|_i \;=\; \sum_{A\ni i}\ \sup_{\eta}|\Phi_A(\eta)| \;<\;\infty
      \qquad(i\in S),\]
    with the locally convex topology generated by the seminorms $(\|\cdot\|_i)_{i\in S}$:
    $\Phi^x\to\Phi$ in $\Bspace$ iff $\|\Phi^x-\Phi\|_i\to0$ for every $i$.
\end{definition}

\begin{proposition}[Georgii (2.11): $\Bspace$ is a Fréchet space]
    \label{prop:exg-b-frechet}
    \uses{def:exg-b-space}
    \lean{Potential.eq_zero_of_forall_seminormAt_eq_zero, Potential.normAt_le_liminf,
      Potential.instT1SpaceSubtypeMemSubmoduleRealAbsolutelySummable,
      Potential.instMetrizableSpaceSubtypeMemSubmoduleRealAbsolutelySummableOfCountable,
      Potential.instCompleteSpaceSubtypeMemSubmoduleRealAbsolutelySummableOfCountable}
    \leanok

    The seminorms $\|\cdot\|_i$ separate points, so $\Bspace$ is Hausdorff; for countable $S$
    it is metrizable and complete.
\end{proposition}
\begin{proof}
    \uses{def:exg-b-space}
    \leanok

    Separation: $\|\Phi\|_i=0$ for all $i$ forces $\Phi_A=0$ for $A\neq\emptyset$ (pick
    $i\in A$), and $\Phi_\emptyset=0$ by definition. Completeness: a Cauchy sequence is Cauchy
    in each $\|\cdot\|_i$, hence termwise Cauchy in $\R$ since $|\Phi_A(\eta)|\le\|\Phi\|_i$
    for $i\in A$. The termwise limit $\Psi$ is absolutely summable because $\|\cdot\|_i$ is
    lower semicontinuous under termwise convergence (Fatou for counting measure) and the
    sequence is $\|\cdot\|_i$-bounded; the same lower semicontinuity applied to
    $\Phi^n-\Phi^m$ gives $\|\Phi^n-\Psi\|_i\le\varepsilon$.
\end{proof}

\subsection*{Georgii (4.23): the Gibbs correspondence $\Gibbs{\cdot}$ on $\Bspace$}

Fix a finite a priori measure $\lambda\in\Meas(E,\mathcal E)$ and $\beta\in\R$; $\gamma^\Phi$
is the Gibbsian specification of $\Phi\in\Bspace$.

\begin{lemma}[Georgii, Comment (4.14)(1)]
    \label{lem:exg-4-14-1}
    \uses{def:exg-b-space, def:exg-finite-volume-distribution}
    \lean{Potential.dominatingMeasure, Potential.locallyEquicontinuous_finiteVolumeDistributions}
    \leanok

    For $\Phi\in\Bspace$ the density of $\gamma^\Phi_\Lambda$ against the free measure is at
    most $e^{2|\beta|\,\|\Phi\|_\Lambda}$, where $\|\Phi\|_\Lambda = \sum_{i\in\Lambda}\|\Phi\|_i$.
    Hence $(\gamma^\Phi_\Lambda(\cdot\mid\eta))_{\Lambda\in\Fin}$ is locally equicontinuous
    for every $\eta$.
\end{lemma}
\begin{proof}
    \leanok

    Georgii (2.14) gives $|H^\Phi_\Lambda|\le\|\Phi\|_\Lambda$, so
    $\rho^\Phi_\Lambda = e^{-\beta H^\Phi_\Lambda}/Z^\Phi_\Lambda
    \le e^{2|\beta|\,\|\Phi\|_\Lambda}$. Apply Georgii (4.12) with $B=\Cfg$.
\end{proof}

\begin{theorem}[Georgii (4.23)(a)]
    \label{thm:exg-4-23-a}
    \uses{thm:exg-4-22, lem:exg-4-14-1, lem:exg-gp-closed, lem:exg-gp-dominated, def:exg-b-space}
    \lean{Potential.GP_gibbsSpecification_nonempty,
      Potential.isCompact_setOf_mem_GP_gibbsSpecification}
    \leanok

    Let $(E,\mathcal E)$ be standard Borel. For every $\Phi\in\Bspace$, $\Gibbs{\Phi}$ is
    non-empty and $\mathcal L$-compact.
\end{theorem}
\begin{proof}
    \uses{thm:exg-4-22, lem:exg-4-14-1, lem:exg-gp-closed, lem:exg-gp-dominated}
    \leanok

    $\gamma^\Phi$ is quasilocal (Georgii, Example (2.25)) and
    $(\gamma^\Phi_\Lambda(\cdot\mid\omega))_\Lambda$ is locally equicontinuous by
    Lemma \ref{lem:exg-4-14-1}, so Theorem \ref{thm:exg-4-22} gives non-emptiness.
    $\Gibbs{\Phi}$ is closed by Lemma \ref{lem:exg-gp-closed} and, by
    Lemmas \ref{lem:exg-gp-dominated} and \ref{lem:exg-4-14-1}, contained in the
    $\mathcal L$-compact set of random fields dominated by the finite measures
    $e^{2|\beta|\,\|\Phi\|_\Lambda}\lambda^S$.
\end{proof}

\begin{theorem}[Georgii (4.23)(b)]
    \label{thm:exg-4-23-b}
    \uses{lem:exg-gp-dominated, lem:exg-4-14-1, def:exg-b-space}
    \lean{Potential.isCompact_closure_iUnion_setOf_mem_GP_of_iSup_normAt_lt_top,
      Potential.isCompact_closure_iUnion_setOf_mem_GP}
    \leanok

    Let $(E,\mathcal E)$ be standard Borel. If $M\subseteq\Bspace$ is bounded, i.e.
    $\sup_{\Phi\in M}\|\Phi\|_i<\infty$ for every $i\in S$, then
    $\Gibbs{M}=\bigcup_{\Phi\in M}\Gibbs{\Phi}$ is relatively $\mathcal L$-compact.
\end{theorem}
\begin{proof}
    \uses{lem:exg-gp-dominated, lem:exg-4-14-1}
    \leanok

    Boundedness gives $B_\Lambda\ge\|\Phi\|_\Lambda$ for all $\Phi\in M$, so by
    Lemmas \ref{lem:exg-4-14-1} and \ref{lem:exg-gp-dominated} the union lies in the
    $\mathcal L$-compact set of random fields dominated by $e^{2|\beta|B_\Lambda}\lambda^S$.
\end{proof}

\begin{theorem}[Georgii (4.23)(c)]
    \label{thm:exg-4-23-c}
    \uses{thm:exg-4-17, def:exg-b-space}
    \lean{Potential.BSpace.gibbsSpecification,
      Potential.BSpace.tendsto_dist_action_gibbsSpecification,
      Potential.BSpace.isClosed_graph_GP}
    \leanok

    With no assumption on $(E,\mathcal E)$, the graph
    $\{(\Phi,\mu):\Phi\in\Bspace,\ \mu\in\Gibbs{\Phi}\}$ is closed in
    $\Bspace\times\Prob(\Cfg,\mathcal F)$.
\end{theorem}
\begin{proof}
    \uses{thm:exg-4-17, lem:exg-mem-gp-iff}
    \leanok

    Let $(\Phi^\alpha,\mu_\alpha)\to(\Phi,\mu)$ in the graph. Index by $\bar D = D\times\Fin$
    with $\Phi^{(\alpha,\Lambda)}=\Phi^\alpha$, $\mu_{(\alpha,\Lambda)}=\mu_\alpha$,
    $\Lambda_{(\alpha,\Lambda)}=\Lambda$; then $\Lambda_{(\alpha,\Lambda)}\to S$, and
    $|H^{\Phi^\alpha}_\Lambda-H^\Phi_\Lambda|\le\|\Phi^\alpha-\Phi\|_\Lambda\to0$ gives
    $\gamma^{\Phi^\alpha}\to\gamma^\Phi$ uniformly on local observables (Georgii (4.19)).
    Since $\mu_\alpha\gamma^{\Phi^\alpha}_\Lambda=\mu_\alpha\to\mu$,
    Theorem \ref{thm:exg-4-17} yields $\mu\in\Gibbs{\Phi}$.
\end{proof}

\begin{theorem}[Georgii (4.23)(d)]
    \label{thm:exg-4-23-d}
    \uses{thm:exg-4-23-c, lem:exg-4-14-1, lem:exg-gp-dominated, prop:exg-b-frechet}
    \lean{Potential.BSpace.isClosed_setOf_exists_mem_GP}
    \leanok

    Let $(E,\mathcal E)$ be standard Borel. For every closed
    $F\subseteq\Prob(\Cfg,\mathcal F)$ the set
    $\Gibbs{\cdot}^{-1}(F)=\{\Phi\in\Bspace:\Gibbs{\Phi}\cap F\neq\emptyset\}$ is closed in
    $\Bspace$; i.e.\ the Gibbs correspondence is upper semicontinuous.
\end{theorem}
\begin{proof}
    \uses{thm:exg-4-23-c, lem:exg-4-14-1, lem:exg-gp-dominated}
    \leanok

    Let $\Phi$ be a cluster point of $\Gibbs{\cdot}^{-1}(F)$ and pick
    $\mu_\Psi\in\Gibbs{\Psi}\cap F$ for each $\Psi$ in that set. Eventually
    $\|\Psi\|_\Lambda\le\|\Phi\|_\Lambda+|\Lambda|$, so by Lemmas \ref{lem:exg-4-14-1} and
    \ref{lem:exg-gp-dominated} the $\mu_\Psi$ are dominated by one family of finite measures,
    hence locally equicontinuous. By Georgii (4.9) they converge along an ultrafilter to some
    $\mu$, which lies in $F$ ($F$ closed) and in $\Gibbs{\Phi}$
    (Theorem \ref{thm:exg-4-23-c}).
\end{proof}

\begin{proposition}[Georgii (4.23) for an arbitrary finite a priori measure]
    \label{prop:exg-4-23-finite-lambda}
    \uses{thm:exg-4-23-a, thm:exg-4-23-b, thm:exg-4-23-c, thm:exg-4-23-d, def:premodif}
    \lean{Specification.lambdaSpecification,
      MeasureTheory.Measure.probNormalize,
      MeasureTheory.Measure.exists_measurable_pos_isProbabilityMeasure_withDensity,
      Potential.gibbsSpecificationOfSigmaFiniteAdmissible,
      Potential.gibbsSpecificationOfFiniteReference,
      Potential.gibbsSpecificationOfFiniteReference_eq_gibbsSpecificationOfAbsolutelySummable,
      Potential.GP_gibbsSpecificationOfFiniteReference_nonempty,
      Potential.isCompact_setOf_mem_GP_gibbsSpecificationOfFiniteReference,
      Potential.BSpace.isClosed_graph_GP_ofFiniteReference,
      Potential.BSpace.isClosed_setOf_exists_mem_GP_ofFiniteReference}
    \leanok

    Theorems \ref{thm:exg-4-23-a}--\ref{thm:exg-4-23-d} hold for every finite non-zero
    $\lambda\in\Meas(E,\mathcal E)$. Finiteness cannot be dropped: by (2.14) an absolutely
    summable potential is $\lambda$-admissible exactly when $\lambda$ is finite.
\end{proposition}
\begin{proof}
    \uses{thm:exg-4-23-a, thm:exg-4-23-b, thm:exg-4-23-c, thm:exg-4-23-d}
    \leanok

    By Georgii, Remark (1.28)(3), replacing $\lambda$ by $r\lambda$ ($r>0$ measurable, finite)
    and $\rho_\Lambda$ by $\rho_\Lambda/\prod_{i\in\Lambda}r(\omega_i)$ leaves the
    $\lambda$-specification unchanged. With $r\equiv\lambda(E)^{-1}$ the normalization absorbs
    the constant, so $\gamma^\Phi$ computed with $\lambda$ equals $\gamma^\Phi$ computed with
    the probability measure $\lambda(E)^{-1}\lambda$, and (a)--(d) transport.
\end{proof}

\subsection*{Georgii Corollary (4.13) in Gibbsian form}

\begin{corollary}[Georgii (4.13), Hamiltonian form]
    \label{cor:exg-4-13}
    \uses{lem:exg-4-14-1, def:exg-b-space}
    \lean{Potential.premodifierNorm_le_of_abs_hamiltonian_le,
      Potential.locallyEquicontinuous_of_confinement_hamiltonian}
    \leanok

    Let $\lambda$ be a probability measure, $(K_\ell)_{\ell\ge1}$ measurable subsets of $E$
    with $\lambda(K_\ell)>0$, $(\Phi^a)_a$ summable $\lambda$-admissible potentials, and
    $\mu_a=\nu_a\gamma^{\Phi^a}_{\Lambda_a}$ with $\Lambda_a\to S$. Suppose
    \begin{enumerate}
        \item $\lim_{\ell\to\infty}\limsup_a\mu_a(\{\omega:\omega_i\notin K_\ell\})=0$ for
              every $i\in S$, and
        \item for every $\Lambda\in\Fin$ there is $\Delta\supseteq\Lambda$ such that for every
              $\ell$ the Hamiltonians $H^{\Phi^a}_\Lambda$ are eventually uniformly bounded on
              $\{\omega:\omega_i\in K_\ell\ \forall i\in\Delta\}$.
    \end{enumerate}
    Then $(\mu_a)_a$ is locally equicontinuous.
\end{corollary}
\begin{proof}
    \leanok

    From (ii), $|H^{\Phi^a}_\Lambda|\le c$ on the box, so the normalized density satisfies
    $\rho^{\Phi^a}_\Lambda\le e^{|\beta|c}\big/\big(e^{-|\beta|c}\lambda(K_\ell)^{|\Lambda|}\big)$
    there: the numerator is at most $e^{|\beta|c}$, and properness of the independent
    specification bounds the partition function on the box below by
    $e^{-|\beta|c}\lambda(K_\ell)^{|\Lambda|}$. This is the density hypothesis of Georgii
    (4.13); (i) is its escape-of-mass hypothesis.
\end{proof}

\subsection*{Georgii Example (4.20): free and periodic boundary conditions}

\begin{example}[Georgii, Example (4.20)(1): free boundary conditions]
    \label{ex:exg-4-20-1}
    \uses{thm:exg-4-17, def:exg-b-space}
    \lean{Potential.truncation, Potential.tail, Potential.abs_hamiltonian_truncation_sub_le,
      Potential.tendsto_tail_atTop, Potential.tendsto_truncationB,
      Potential.tendsto_dist_action_truncation,
      Potential.mem_GP_of_mapClusterPt_truncation,
      Potential.mem_GP_of_mapClusterPt_truncation_finiteVolumeDistributions}
    \leanok

    For $\Phi\in\Bspace$ and $\Delta\in\Fin$ let $\Phi^\Delta_A=\Phi_A$ if $A\subseteq\Delta$
    and $\Phi^\Delta_A=0$ otherwise. Then
    \[\big|H^{\Phi^\Delta}_\Lambda-H^\Phi_\Lambda\big|\ \le\
      \sum_{A\cap\Lambda\neq\emptyset,\ A\not\subseteq\Delta}\ \sup_\eta|\Phi_A(\eta)|
      \ \xrightarrow[\ \Delta\uparrow S\ ]{}\ 0,\]
    so $\Phi^\Delta\to\Phi$ in $\Bspace$ and $\gamma^{\Phi^\Delta}\to\gamma^\Phi$ uniformly
    on local observables. Every cluster point of the free-boundary net
    $\Delta\mapsto\nu_\Delta\gamma^{\Phi^\Delta}_\Delta$, in particular of
    $(\gamma^{\Phi^\Delta}_\Delta(\cdot\mid\omega))_\Delta$, belongs to $\Gibbs{\Phi}$.
\end{example}
\begin{proof}
    \uses{thm:exg-4-17}
    \leanok

    The Hamiltonians differ only in the terms $\Phi_A$ with $A\cap\Lambda\neq\emptyset$,
    $A\not\subseteq\Delta$, and the tail of the absolutely convergent interaction series over
    $\{A:A\not\subseteq\Delta\}$ tends to $0$ as $\Delta\uparrow S$. Georgii (4.19) gives
    uniform convergence $\gamma^{\Phi^\Delta}\to\gamma^\Phi$, and the cluster-point form of
    Theorem \ref{thm:exg-4-17} concludes.
\end{proof}

\begin{corollary}[Free-boundary existence]
    \label{cor:exg-free-boundary-existence}
    \uses{ex:exg-4-20-1, thm:exg-4-22, lem:exg-4-14-1}
    \lean{Potential.locallyEquicontinuous_truncation_bindPM,
      Potential.exists_mem_GP_mapClusterPt_truncation}
    \leanok

    Let $(E,\mathcal E)$ be standard Borel and $\Phi\in\Bspace$. For any boundary fields
    $(\nu_\Delta)_\Delta$ the free-boundary net $\Delta\mapsto\nu_\Delta\gamma^{\Phi^\Delta}_\Delta$
    is locally equicontinuous, has a cluster point, and every cluster point lies in
    $\Gibbs{\Phi}$.
\end{corollary}
\begin{proof}
    \uses{ex:exg-4-20-1, thm:exg-4-22, lem:exg-4-14-1}
    \leanok

    $\|\Phi^\Delta\|_i\le\|\Phi\|_i$, so on $\mathcal F_\Lambda$-events every
    $\gamma^{\Phi^\Delta}_\Lambda$ is dominated by $e^{2|\beta|\|\Phi\|_\Lambda}\lambda^S$, and
    by consistency so is $\nu_\Delta\gamma^{\Phi^\Delta}_\Delta$ once $\Lambda\subseteq\Delta$.
    Theorem \ref{thm:exg-4-22} applies with the uniform convergence of
    Example \ref{ex:exg-4-20-1}.
\end{proof}

\begin{example}[Georgii, Example (4.20)(2): periodic boundary conditions]
    \label{ex:exg-4-20-2}
    \uses{thm:exg-4-17, ex:exg-4-20-1, def:exg-b-space}
    \lean{Potential.IsTorusReduction, Potential.IsAnchor, Potential.periodicExtend,
      Potential.periodicModification, Potential.normAt_periodicModification_le,
      Potential.isAbsolutelySummable_periodicModification,
      Potential.abs_hamiltonian_periodicModification_sub_le,
      Potential.mem_GP_of_mapClusterPt_periodicModification,
      Potential.mem_GP_of_mapClusterPt_latticePeriodic_finiteVolumeDistributions}
    \leanok

    Let $\Phi\in\Bspace$ be shift invariant and $\Delta$ a box identified with the torus
    $S/p\!\cdot\!S$ through a reduction $\pi$ with periodic continuation
    $\tilde\sigma_\Delta(\omega)_i=\omega_{\pi(i)}$, the representatives $\mathcal R(A)$ of
    the translation classes being chosen by a translation-equivariant anchor. The
    $\Delta$-periodic modification
    $\tilde\Phi^\Delta_A=\sum_{B\in\mathcal R(A)}\Phi_B\circ\tilde\sigma_\Delta$ satisfies
    $\|\tilde\Phi^\Delta\|_i\le\|\Phi\|_i$, hence lies in $\Bspace$, and
    \[\big|H^{\tilde\Phi^\Delta}_\Lambda-H^\Phi_\Lambda\big|\ \le\ 2
      \sum_{A\cap\Lambda\neq\emptyset,\ A\not\subseteq\Delta}\ \sup_\eta|\Phi_A(\eta)| .\]
    Every cluster point of $\Delta\mapsto\nu_\Delta\gamma^{\tilde\Phi^\Delta}_\Delta$,
    $\Delta\uparrow S$, belongs to $\Gibbs{\Phi}$. On $S=\mathbb Z^d$ this applies to the
    cubes $\Delta_n=[-(n+1),n+1)^d$ with coordinatewise reduction and lexicographic anchor,
    with any boundary condition $\omega$.
\end{example}
\begin{proof}
    \uses{thm:exg-4-17, ex:exg-4-20-1}
    \leanok

    Reindexing along the anchor gives $\|\tilde\Phi^\Delta\|_i\le\|\Phi\|_i$. In the
    difference of Hamiltonians the terms with $A\subseteq\Delta$ cancel against the
    representative $B=A$, leaving two sums each dominated by the tail of
    Example \ref{ex:exg-4-20-1}. Georgii (4.19) gives uniform convergence and the
    cluster-point form of Theorem \ref{thm:exg-4-17} concludes. The cubes $\Delta_n$ are
    cofinal in Georgii's directed set of rectangular boxes.
\end{proof}

\begin{remark}[Georgii (4.20): the boundary condition drops out]
    \label{rmk:exg-4-20-supported}
    \uses{ex:exg-4-20-1, ex:exg-4-20-2}
    \lean{Potential.gibbsSpecificationOfAbsolutelySummable_apply_eq_of_supported,
      Potential.gibbsSpecification_truncation_apply_eq,
      Potential.gibbsSpecification_periodicModification_apply_eq}
    \leanok

    For a potential $\Psi$ supported in $\Delta$, in particular $\Phi^\Delta$ and
    $\tilde\Phi^\Delta$, the restriction of $\gamma^\Psi_\Delta(\cdot\mid\omega)$ to
    $\mathcal F_\Delta$ does not depend on $\omega$.
\end{remark}

\subsection*{Per-site a priori measures}

\begin{definition}[Inhomogeneous independent specification]
    \label{def:exg-isssd-family}
    \uses{def:isssd, def:isssd-spec, def:juxtaposition, def:markov-spec}
    \lean{Specification.isssdFamilyFun, Specification.isssdFamily,
      Specification.isssdFamily_const, Specification.hasFreeMeasure_isssdFamily}
    \leanok

    For a family $(\lambda_i)_{i\in S}$ of single-spin probability measures, the kernel
    resampling the sites of $\Lambda$ independently, site $i$ from $\lambda_i$, and keeping
    the boundary condition outside $\Lambda$, is a specification. A constant family gives
    $\ISSSD(\lambda)$; on $\mathcal F_\Lambda$-events the kernel equals
    $\bigotimes_i\lambda_i$ regardless of the boundary condition.
\end{definition}

\begin{theorem}[Georgii (4.23)(a) over per-site a priori measures]
    \label{thm:exg-4-23-a-family}
    \uses{def:exg-isssd-family, thm:exg-4-22, thm:exg-4-23-a, lem:exg-gp-dominated}
    \lean{Potential.gibbsSpecificationFamily,
      Potential.isQuasilocal_gibbsSpecificationFamily,
      Potential.GP_gibbsSpecificationFamily_nonempty,
      Potential.isCompact_setOf_mem_GP_gibbsSpecificationFamily}
    \leanok

    Let $(E,\mathcal E)$ be standard Borel. The Gibbsian specification of $\Phi\in\Bspace$
    relative to a family $(\lambda_i)_{i\in S}$ of single-spin probability measures is
    quasilocal, and its set of Gibbs measures is non-empty and $\mathcal L$-compact.
\end{theorem}
\begin{proof}
    \uses{def:exg-isssd-family, thm:exg-4-22, lem:exg-gp-dominated}
    \leanok

    The proof of Theorem \ref{thm:exg-4-23-a} uses two properties of the reference
    specification: it resamples volumes, giving quasilocality of the modification, and it
    forgets the boundary condition on $\mathcal F_\Lambda$-events, giving the density bound
    $e^{2|\beta|\|\Phi\|_\Lambda}$ against $\bigotimes_i\lambda_i$ and hence local
    equicontinuity. Both hold for Definition \ref{def:exg-isssd-family}, so
    Theorem \ref{thm:exg-4-22} and Lemma \ref{lem:exg-gp-dominated} apply.
\end{proof}

\begin{corollary}[Absorbing the self-energies: the existence half of Georgii (8.39)]
    \label{cor:exg-self-energy}
    \uses{thm:exg-4-23-a-family, def:exg-isssd-family, def:premodif}
    \lean{Potential.selfEnergyWeight, Potential.selfEnergyMeasure,
      Potential.lambdaSpecification_eq_gibbsSpecificationFamily,
      Potential.lintegral_selfEnergyWeight_ne_top,
      Potential.GP_lambdaSpecification_nonempty}
    \leanok

    Let $(E,\mathcal E)$ be standard Borel, $\lambda$ $\sigma$-finite and non-zero, $\Phi$
    summable and $\lambda$-admissible, and $\eta_0\in\Cfg$. If the recentred many-body part
    $(\Phi^{\mathrm{many}})^{\eta_0}$ of $\Phi$ is absolutely summable, the
    $\lambda$-specification of $\Phi$ has a Gibbs measure. On $\mathbb Z$ or $\mathbb N$,
    condition (8.40) of Theorem \ref{thm:onedim-839} implies that summability.
\end{corollary}
\begin{proof}
    \uses{thm:exg-4-23-a-family}
    \leanok

    Split $\Phi$ into its one-body part $(\Phi_{\{i\}})_i$ and its many-body part, recentred
    at $\eta_0$. The Boltzmann factor is a constant times the Boltzmann factor of the recentred
    many-body part times $\prod_{i\in\Lambda}e^{-\beta\Phi_{\{i\}}}$; by Remark (1.28)(3) the
    last factor is the weight by which the reference kernel changes when $\lambda$ is replaced
    at site $i$ by $c_i^{-1}e^{-\beta\Phi_{\{i\}}}\lambda$. The constants $c_i$ are finite
    because $\lambda$-admissibility at single sites bounds
    $\lambda(e^{-\beta\Phi_{\{i\}}})$. Hence the $\lambda$-specification of $\Phi$ is the
    Gibbsian specification of the recentred many-body part over the per-site measures, and
    Theorem \ref{thm:exg-4-23-a-family} applies.
\end{proof}

\subsection*{Two remarks}

\begin{remark}[The weak topology is kept separate]
    \label{rmk:exg-weak-topology}
    \uses{lem:exg-mem-gp-iff}
    \lean{Specification.continuous_bindPM,
      MeasureTheory.GibbsMeasure.isClosed_GP,
      MeasureTheory.GibbsMeasure.isCompact_GP,
      MeasureTheory.GibbsMeasure.existence_of_gibbsMeasure_compact_weak,
      MeasureTheory.GibbsMeasure.existence_of_gibbsMeasure_of_isTight_weak}
    \leanok

    If $\gamma$ is Feller, $\mu\mapsto\mu\gamma_\Lambda$ is weakly continuous, so
    $\Gibbs{\gamma}$ is weakly closed, and weak compactness (from compactness of $E$, or
    tightness of the finite-volume distributions) yields a Gibbs measure. Weak compactness is
    distinct from the $\mathcal L$-compactness of Theorem \ref{thm:exg-4-22}.
\end{remark}

\begin{remark}[Georgii's Appendix 4.A]
    \label{rmk:exg-appendix-4a}

    The results of Appendix 4.A --- the isomorphism theorem (4.A6), countable cores (4.A9),
    Theorem (4.A11), Corollary (4.A13) --- enter in Chapters 7, 14 and 15.
\end{remark}
